% chapters/02-debugging.tex -- part A. 디버깅: 방법, 자기 proxy 사건, 열세 사례.
\chapter{디버깅}\label{ch:debugging}

% src: books/debugging/README.md (Result line, Index)
\paragraph{결과.}
2026/10/13에 막힌 곳이 13건이었다.
11건은 고쳤고 1건은 열려 있다.
열린 것은 사례 3의 내용 문제다(CDN의 썸네일이 없다).
사례 13은 고장이 아니었다. 옛 오류 줄을 새 오류로 읽은 것이어서 고칠 것이 없었다.

\section{방법: 추측하기 전에 돌리는 순서}\label{sec:dbg-method}

\begin{inbrief}
  \item[Why] 추측으로 고치면 같은 고장을 다시 만난다.
  \item[When] 증상이 보이고 원인을 모를 때다.
  \item[Where] 앱, VM, 터널, Clerk, PM2 어디서든 쓴다.
  \item[How] 증상을 줄이고, 길을 나누고, 오류를 글자 그대로 읽고, 하나만 바꾸고, 실행 중인 값을 확인한다.
\end{inbrief}
% src: books/debugging/00-method.md "The loop"

\Cref{fig:dbg-loop}은 방법을 한 줄로 보인다.
일곱 습관이 모두 이 줄 위에 있다.

\begin{figure}[H]
  \centering
  \begin{tikzpicture}[node distance=0.3cm and 0.3cm, font=\footnotesize]
    \node[flowbox, minimum width=1.9cm] (s1) {증상};
    \node[flowbox, minimum width=1.9cm, right=of s1] (s2) {줄이기\\\scriptsize\texttt{curl -m}};
    \node[flowbox, minimum width=1.9cm, right=of s2] (s3) {길 나누기\\\scriptsize middleware 유무\\\scriptsize VM 대 터널};
    \node[flowbox, minimum width=1.9cm, right=of s3] (s4) {오류를 글자 그대로\\\scriptsize host, scheme, port};
    \node[flownew, minimum width=1.9cm, right=of s4] (s5) {하나만 바꾼다};
    \node[flownew, minimum width=1.9cm, right=of s5] (s6) {실행 중인 값 확인\\\scriptsize\texttt{pm2 env}, \texttt{ss}};
    \foreach \a/\b in {s1/s2,s2/s3,s3/s4,s4/s5,s5/s6} \draw[flowarrow] (\a) -- (\b);
    \node[flowalt, below=0.7cm of s4, minimum width=4.2cm] (src) {로그와 문서가 어긋나면\\프레임워크 소스를 읽는다};
    \draw[flowarrow, dashed] (s4) -- (src);
  \end{tikzpicture}
  \flowlegend
  \caption{증상에서 확인까지의 디버깅 순서. 파란 상자는 고치는 단계, 점선 상자는 막다른 길에서 가는 곳이다.}
  \label{fig:dbg-loop}
  % src: books/debugging/00-method.md "The loop"
\end{figure}

\subsection{일곱 습관}\label{sec:dbg-habits}

\paragraph{결과.}
일곱 습관은 모두 이 책의 사례 중 최소 하나를 줄였다.
% src: books/debugging/00-method.md "Result"

\begin{table}[H]
  \caption{일곱 습관, 하는 법, 사례에서 나온 실제 값. 사례 번호는 \Cref{tab:cases}의 번호다.}
  \label{tab:dbg-habits}
  \centering\footnotesize
  \begin{tabularx}{\linewidth}{@{}L{0.5cm}L{3.2cm}L{3.9cm}Y L{1.1cm}@{}}
    \toprule
    \# & 습관 & 하는 법 & 사례의 실제 값 & 사례 \\
    \midrule
    1 & 모든 \code{curl}에 시간 제한 \code{-m}을 쓴다 & \code{curl -m 20 http://127.0.0.1:3000/} & 20 s 뒤 \code{000}(멈춤). \code{-m}이 없으면 터미널이 끝없이 기다린다 & 1 \\
    \midrule
    2 & middleware를 지나는 경로와 안 지나는 경로를 비교한다 & 지나는 \code{/}, 안 지나는 \code{/robots.txt} & \code{/}는 멈췄고 \code{/robots.txt}는 0.25 s에 \code{200}. 결함은 middleware에 있다 & 1 \\
    \midrule
    3 & 먼저 VM 안에서, 그 다음 터널로 시험한다 & VM에서 \code{curl http://127.0.0.1:3000/}, 그 뒤 Mac에서 공개 host & VM에서 실패하면 앱 결함이다. Mac에서만 새 실패가 나오면 터널을 본다 & 1 \\
    \midrule
    4 & 오류의 URL을 글자 그대로 읽는다 & 오류 문구에서 host, scheme, port를 옮겨 적고 각 부분을 누가 만들었나 묻는다 & \code{Failed to proxy https://localhost:3000/}: \code{https}는 \code{X-Forwarded-Proto}에서 왔고 \code{localhost}는 우리가 정한 \code{HOSTNAME}이 아니다. \code{https://local.prettylog.com:3000}의 포트를 Clerk가 거절했다 & 1, 5 \\
    \midrule
    5 & 로그와 문서가 어긋나면 프레임워크 소스를 읽는다 & stack이 가리키는 \code{node_modules} 파일과 줄을 연다 & \code{resolve-routes.js:117}과 \code{next-server.js:1132}로 정리됐다. 어떤 문서에도 없었다 & 1 \\
    \midrule
    6 & 한 번에 하나만 바꾸고 다시 확인한다 & 값 하나를 바꾸고 같은 검사를 다시 돌린다 & 첫 가설(cloudflared의 Host 헤더)은 헤더 없는 \code{curl}도 실패해서 버렸다. 되돌린 시도(\code{unstable_rethrow}, 사례 3)도 결과로 적었다 & 1, 7, 8 \\
    \midrule
    7 & 파일이 아니라 실행 중인 값을 확인한다 & \code{pm2 env}, \code{ss -ltnp}로 프로세스의 실제 값을 본다 & 디스크의 파일은 프로세스 환경이 아니다 & 8, 9 \\
    \bottomrule
  \end{tabularx}
  % src: books/debugging/00-method.md "Table", "How to use each habit"
\end{table}

\subsection{증상 읽는 법}\label{sec:dbg-reading}

\begin{table}[H]
  \caption{증상 한 줄과 먼저 의심할 곳.}
  \label{tab:dbg-reading}
  \centering\small
  \begin{tabularx}{\linewidth}{@{}L{5.2cm}Y@{}}
    \toprule
    보이는 것 & 의심할 곳 \\
    \midrule
    시간 제한 뒤 \code{000} & 멈춤. proxy 순환이나 막힌 호출을 찾는다 \\
    \midrule
    바로 \code{000} & 거절됨. 그 주소에서 받는 것이 없다 \\
    \midrule
    정적 파일은 \code{200}, \code{/}는 멈춤 & middleware \\
    \midrule
    CORS 헤더 없는 \code{500} & 실행 중인 목록에 origin이 없다 \\
    \midrule
    \code{401} & Clerk 인스턴스가 다르거나 \code{azp}가 목록에 없다 \\
    \midrule
    release 뒤 오류 로그의 \code{ETIMEDOUT} & 옛 줄일 수 있다. 파일 수정 시각과 줄 수를 먼저 본다 (\Cref{sec:dbg-case13}) \\
    \bottomrule
  \end{tabularx}
  % src: books/debugging/00-method.md "Reading"
\end{table}

\section{깊이 보기: blog-web이 자기 자신에게 proxy한 사건}\label{sec:dbg-selfproxy}

% src: deploy/gcp/web-self-proxy.md Status; books/debugging/01-web-self-proxy.md Result
\paragraph{결과.}
2026/10/13에 고쳤다.
\texttt{ecosystem.config.cjs}의 \texttt{HOSTNAME}을 \texttt{'127.0.0.1'}에서 \texttt{'localhost'}로 바꿨다(\commit{3a6463a}).
VM과 바깥에서 모두 확인했다.

\subsection{증상}\label{sec:dbg-sp-symptom}

\begin{table}[H]
  \caption{blog-web 첫 배포의 증상. 열은 어디서 보였나, 무엇이 보였나이다.}
  \label{tab:sp-symptom}
  \centering\small
  \begin{tabularx}{\linewidth}{@{}L{5.0cm}Y@{}}
    \toprule
    어디서 & 무엇이 \\
    \midrule
    \code{https://my.prettylog.com} (첫 배포) & \code{502} 또는 응답 없음 \\
    \midrule
    VM의 \code{curl -m 20 http://127.0.0.1:3000/} & 20 s 뒤 \code{000}(멈춤), 또는 바로 \code{000}(거절), 한 번은 \code{500} \\
    \midrule
    VM의 \code{curl http://127.0.0.1:3000/robots.txt} & 0.25 s에 \code{200}. 그 경로는 middleware를 지나지 않는다 \\
    \midrule
    터널 경유 \code{pm2 logs blog-web --err} & \code{Failed to proxy https://localhost:3000/ Error: write EPROTO ... wrong version number} \\
    \midrule
    로컬 curl 경유 \code{pm2 logs blog-web --err} & \code{Failed to proxy http://localhost:3000/ Error: connect ECONNREFUSED 127.0.0.1:3000} \\
    \midrule
    Express(\code{blog-api}), 메모리 & 정상. \code{/ping}이 공개와 로컬 모두 \code{200}, 475 MiB available \\
    \bottomrule
  \end{tabularx}
  % src: deploy/gcp/web-self-proxy.md "Symptom"
\end{table}

\paragraph{읽는 법.}
\texttt{/}만 멈추고 \texttt{/robots.txt}는 빠르다.
경로에 따라 갈리면 그 경로만 지나는 층을 본다.
여기서는 middleware였다.
% src: books/debugging/01-web-self-proxy.md Lesson

\subsection{원인}\label{sec:dbg-sp-cause}

\Cref{fig:sp-cause}는 요청이 자기 자신으로 되돌아가는 길을 보인다.
핵심은 두 origin이 서로 달랐다는 점이다.

\begin{figure}[H]
  \centering
  \begin{tikzpicture}[font=\small, node distance=0.6cm]
    \node[flowbox, minimum width=11cm] (req) {요청이 Next router에 들어온다};
    \node[flowbox, minimum width=11cm, below=of req] (own) {Next router의 자기 origin = \texttt{<proto>://127.0.0.1:3000}\\\scriptsize(\texttt{HOSTNAME=127.0.0.1}에서 만든다)};
    \node[flowbox, minimum width=11cm, below=of own] (mw) {middleware가 보는 요청 URL = \texttt{<proto>://localhost:3000}\\\scriptsize(Next의 기본 host)};
    \node[flowbox, minimum width=11cm, below=of mw] (clerk) {Clerk \texttt{clerkMiddleware}는 \texttt{NextResponse.next()}마다\\\texttt{x-middleware-rewrite} = 요청 URL로 둔다};
    \node[badbox, minimum width=11cm, below=of clerk] (ext) {router: rewrite origin(\texttt{localhost}) $\neq$ 자기 origin(\texttt{127.0.0.1})\\$\to$ ``외부 rewrite''로 보고 \texttt{<proto>://localhost:3000}으로 HTTP proxy};
    \node[badbox, minimum width=5.2cm, below=0.7cm of ext.south west, anchor=north west] (e1) {\texttt{<proto>} = https\\\texttt{EPROTO}\\\scriptsize \texttt{X-Forwarded-Proto}가\\cloudflared에서 온다};
    \node[badbox, minimum width=5.2cm, below=0.7cm of ext.south east, anchor=north east] (e2) {\texttt{<proto>} = http\\\texttt{ECONNREFUSED} 또는 멈춤};
    \foreach \a/\b in {req/own,own/mw,mw/clerk,clerk/ext} \draw[flowarrow] (\a) -- (\b);
    \draw[flowarrow] (ext.south -| e1) -- (e1);
    \draw[flowarrow] (ext.south -| e2) -- (e2);
  \end{tikzpicture}
  \caption{blog-web이 자기 자신에게 proxy한 원인. 같은 서버를 가리키는 두 이름(\texttt{127.0.0.1}, \texttt{localhost})이 다른 origin으로 취급됐다.}
  \label{fig:sp-cause}
  % src: deploy/gcp/web-self-proxy.md "Cause"
\end{figure}

\paragraph{근거: 소스의 어느 줄인가.}
아래 줄은 Next 16.2.9와 \texttt{@clerk/nextjs} 7.5.9의 것이다.
% src: deploy/gcp/web-self-proxy.md "Cause" table

\begin{table}[H]
  \caption{원인을 받치는 소스 줄. 파일 경로와 줄 번호는 원문 그대로다.}
  \label{tab:sp-source}
  \centering\footnotesize
  \begin{tabularx}{\linewidth}{@{}YL{6.2cm}@{}}
    \toprule
    단계 & 소스 \\
    \midrule
    router origin = \code{protocol://} + \code{formatHostname(opts.hostname)} + port & \code{next/dist/server/lib/router-utils/resolve-routes.js:117} \\
    \midrule
    \code{X-Forwarded-Proto}에 \code{https}가 있으면 \code{protocol} = \code{https} & 같은 파일, 115번 줄 \\
    \midrule
    middleware URL = \code{initProtocol://} + \code{(this.fetchHostname || 'localhost')} + port & \code{next/dist/server/next-server.js:1132} \\
    \midrule
    origin이 다르면 절대 URL로 남아 proxy된다 & \code{next/dist/shared/lib/router/utils/relativize-url.js} (\code{parseRelativeURL}) \\
    \midrule
    Clerk가 \code{next()}를 같은 URL의 rewrite로 바꾼다 & \code{@clerk/nextjs/dist/esm/server/utils.js:44-67} (\code{decorateRequest}) \\
    \bottomrule
  \end{tabularx}
  % src: deploy/gcp/web-self-proxy.md "Cause" table
\end{table}

\begin{callout}[증명하지 못한 것]
standalone 서버에서 \texttt{fetchHostname}이 비어 있다는 점은 증명하지 못했다.
증거는 로그다. proxy된 URL에 \texttt{localhost}가 찍혔고 \texttt{HOSTNAME}은 \texttt{127.0.0.1}이었다.
% src: deploy/gcp/web-self-proxy.md "Not proven"
\end{callout}

\paragraph{왜 \texttt{/robots.txt}는 멀쩡했나.}
middleware matcher(\texttt{proxy.ts})에 든 경로만 이 길을 탄다.
\texttt{/robots.txt}, \texttt{/\_next/*}, 확장자가 있는 파일은 matcher 밖이다.
% src: deploy/gcp/web-self-proxy.md "Cause" last paragraph

\paragraph{버린 첫 가설.}
cloudflared의 Host 헤더가 원인이라고 먼저 짐작했다.
헤더 없는 평범한 \texttt{curl}도 실패해서 버렸다.
% src: books/debugging/01-web-self-proxy.md "Cause"

\subsection{고침}\label{sec:dbg-sp-fix}

\paragraph{무엇을 바꿨나.}
\texttt{deploy/gcp/ecosystem.config.cjs}의 앱 \texttt{blog-web} 한 줄이다.
% src: deploy/gcp/web-self-proxy.md "Fix"

\begin{lstlisting}[language={}]
-                HOSTNAME: '127.0.0.1',
+                HOSTNAME: 'localhost',
\end{lstlisting}

\paragraph{왜 통하나.}
두 origin이 모두 \texttt{<proto>://localhost:3000}이 된다.
그러면 rewrite가 내부에 남는다.
\texttt{localhost}는 VM에서 \texttt{127.0.0.1}로 풀리므로 서버는 여전히 loopback에서만 받는다.
% src: deploy/gcp/web-self-proxy.md "Fix"

\paragraph{어떻게 적용했나.}
다시 빌드하지 않는다. 설정 파일을 올리고 PM2를 다시 시작한다.
% src: deploy/gcp/web-self-proxy.md "Fix" Apply

\begin{lstlisting}[language={}]
# Mac, repo root
gcloud compute scp deploy/gcp/ecosystem.config.cjs $VM:~/prettylog/ "${Z[@]}"
# VM
pm2 restart ~/prettylog/ecosystem.config.cjs --only blog-web --update-env
pm2 env $(pm2 id blog-web | tr -dc '0-9') | grep HOSTNAME     # HOSTNAME: localhost
getent ahosts localhost | head -2                             # 127.0.0.1 first
ss -ltnp | grep ':3000'                                       # 127.0.0.1:3000, not 0.0.0.0
\end{lstlisting}

\paragraph{버린 대안.}
\Cref{tab:sp-rejected}가 세 가지를 보인다.

\begin{table}[H]
  \caption{버린 대안과 이유.}
  \label{tab:sp-rejected}
  \centering\small
  \begin{tabularx}{\linewidth}{@{}L{4.6cm}Y@{}}
    \toprule
    대안 & 버린 이유 \\
    \midrule
    \code{HOSTNAME=0.0.0.0} & Next가 이것도 \code{localhost}로 옮기지만 모든 interface에서 받는다. \code{:3000}을 GCP 방화벽만 지키게 된다 \\
    \midrule
    \code{next.config.ts}의 \code{skipProxyUrlNormalize} & 다시 빌드해야 하고 middleware가 URL을 보는 방식이 바뀐다. 대비책으로 남겼다 \\
    \midrule
    Cloudflare tunnel의 ``HTTP Host Header'' & Host 헤더는 두 origin 어디에도 쓰이지 않아 도움이 안 된다 \\
    \bottomrule
  \end{tabularx}
  % src: deploy/gcp/web-self-proxy.md "Rejected"
\end{table}

\subsection{확인}\label{sec:dbg-sp-verify}

\begin{table}[H]
  \caption{고친 뒤 확인(2026/10/13). 기대는 모두 실제 값과 맞았다.}
  \label{tab:sp-verify}
  \centering\footnotesize
  \begin{tabularx}{\linewidth}{@{}L{6.4cm}Y@{}}
    \toprule
    검사 & 결과 \\
    \midrule
    \code{pm2 env ... | grep HOSTNAME} & \code{HOSTNAME: localhost} \\
    \midrule
    \code{getent ahosts localhost} & \code{127.0.0.1}이 먼저 \\
    \midrule
    \code{ss -ltnp | grep :3000} & \code{127.0.0.1:3000} (loopback만) \\
    \midrule
    VM에서 \code{curl http://127.0.0.1:3000/} 3번 & \code{200}, 1.49 s, 0.51 s, 0.35 s \\
    \midrule
    \code{pm2 logs blog-web --err | grep -c 'Failed to proxy'} & \code{0} \\
    \midrule
    Mac에서 \code{https://my.prettylog.com}의 \code{/}, \code{/blog}, 글 하나, \code{/about} & 모두 \code{200} (0.7--1.9 s) \\
    \bottomrule
  \end{tabularx}
  % src: deploy/gcp/web-self-proxy.md "Verified (2026-10-13)"
\end{table}

\paragraph{덧붙임.}
바꾸는 동안 바깥에서 \texttt{500 Internal Server Error}가 한 번 보였다.
옛 프로세스의 응답이었고, 다시 시작한 뒤에는 사라졌다.
\texttt{pm2 save}로 새 프로세스 목록을 저장했다.
% src: deploy/gcp/web-self-proxy.md "Verified" last paragraph

\section{열세 사례}\label{sec:dbg-cases}

\paragraph{결과.}
13건 중 11건은 고쳤고 1건(사례 3)은 내용 문제가 열려 있다.
사례 13은 조치가 필요 없었다.
\Cref{tab:cases}의 마지막 열에 ``증명하지 못한 것''을 소스의 표식 그대로 적었다.
소스가 ``브리프에서 받음, 확인 못 함''이라 한 줄은 고쳤다고 올려 쓰지 않았다.
% src: books/debugging/README.md "Result", "How to read a chapter" last line

\begin{longtable}{@{}L{0.5cm}L{3.0cm}L{3.6cm}L{3.3cm}L{3.4cm}@{}}
  \caption{열세 사례: 증상, 원인, 고침, 확인 상태. 에러 문구, 경로, 식별자는 원문 그대로다.}\label{tab:cases}\\
  \toprule
  \# & 증상 & 원인 & 고침 & 확인 상태 \\
  \midrule
  \endfirsthead
  \toprule
  \# & 증상 & 원인 & 고침 & 확인 상태 \\
  \midrule
  \endhead
  \bottomrule
  \endfoot
  % src: books/debugging/README.md "Index" and chapters 01-13 "Evidence"
  1 & 페이지가 20 s 멈춤. \code{/robots.txt}는 200 & Next router origin \code{127.0.0.1}과 middleware URL \code{localhost}가 달라 Clerk의 rewrite가 외부로 보인다 & \commit{3a6463a}, \code{HOSTNAME: 'localhost'} & 고침. VM 3회 200, \code{Failed to proxy} 0. 증명 안 됨: standalone에서 \code{fetchHostname}이 비어 있는지 \\
  \midrule
  2 & API가 죽으면 홈 전체가 오류 페이지로 감 & \code{getFeed}가 가드 없이 \code{TypeError: fetch failed}를 던짐 & \commit{77315eb} & 고침. 근거는 소스 줄 27--45. 실행 기록은 소스에 없다 \\
  \midrule
  3 & 글 페이지가 오류 페이지로 감 & 상대 이미지 참조가 글 route에 걸려 API가 404, 코드가 일반 Error를 다시 던짐 & \commit{42ba55c} (페이지) & 페이지는 고침(dev 서버 측정: status 200, \code{noindex}). \textbf{내용 문제는 열림}: 오너가 A, B, C를 아직 안 골랐다 \\
  \midrule
  4 & 로컬 web이 prod API에서 401 & 토큰이 Clerk dev 인스턴스 것인데 API는 live key로 검증 & \code{LOCAL_PROD_API.md}의 mode C (\commit{cc5aead}) & 고침. 토큰 \code{iss}, \code{azp}를 읽었고 로그인, 댓글 쓰기, CORS preflight 204가 끝까지 동작(2026/10/13) \\
  \midrule
  5 & Clerk: ``Production Keys are only allowed for domain'' & origin에 포트 \code{:3000}이 붙음 & 443에서 서빙 & 고침. 근거는 \code{LOCAL_PROD_API.md}의 표 행(2026/10/13) \\
  \midrule
  6 & \code{listen EACCES ... 127.0.0.1:443} & 이 Mac에서 일반 사용자는 \code{0.0.0.0:443}만 bind할 수 있음 & \code{-H 0.0.0.0} & 고침. \textbf{macOS가 왜 다른지는 확인 안 함} \\
  \midrule
  7 & 인증서 경고, 그 뒤 HMR 차단 & Next 인증서에 host 없음, \texttt{allowed\-DevOrigins}에 host 없음 & mkcert 인증서, \texttt{allowed\-DevOrigins} & 고침. \code{openssl x509 ... -ext subjectAltName}에 \code{DNS:local.prettylog.com}이 나와야 한다. 인증서 만료 2029/01/13 \\
  \midrule
  8 & preflight가 500, CORS 헤더 없음 & 실행 중인 \code{ALLOWED_ORIGINS}에 origin이 없음. restart가 bare 이름이라 아무것도 다시 시작하지 않음 & 전체 ecosystem 경로 + \code{--update-env} & 고침. preflight 204. \textbf{``500 대신 403'' 제안은 브리프 출처, 이 사례에서는 안 함.} 뒤에 항목 G가 403으로 바꿨다(\commit{4d137f2}) \\
  \midrule
  9 & \code{blog-api}가 계속 재시작(\code{pm2 ls}의 재시작 횟수 15) & \code{api.env}의 숫자 키에 숫자가 아닌 값 & 숫자만 남김 & 고침. \textbf{재시작 횟수 15는 브리프 출처, 실제 잘못된 값도 확인 못 함} \\
  \midrule
  10 & \code{gcloud}가 문자열 전체를 zone으로 받음 & zsh는 따옴표 없는 문자열 변수를 나누지 않음 & 배열 \code{Z=(...)} & 고침. 근거는 \code{deploy/gcp/README.md} 57번 줄. \textbf{gcloud 오류 문구는 소스에 없어 확인 못 함} \\
  \midrule
  11 & setup-gcloud: ``problem refreshing your current auth tokens'' & STS API와 IAM Credentials API가 꺼져 있었음 & \code{gcloud services enable ...} & 고침. Actions run 38354550435가 issue \#3을 열었다(4.5 MB / 7일, 0 ERROR, OK). \textbf{issue \#2의 상세는 브리프 출처, 확인 못 함} \\
  \midrule
  12 & \code{ERR_PNPM_RECURSIVE_EXEC_FIRST_FAIL} & \code{next}가 repo root에 설치돼 있지 않음 & \code{pnpm --filter ... exec} & 고침. 근거는 \code{LOCAL_PROD_API.md} 4단계 \\
  \midrule
  % src: scratchpad web-release-2026-10-13.facts.md §4
  13 & release 뒤 \code{pm2 logs blog-web --err}에 \code{ETIMEDOUT} & 고장이 아님. 로그 724줄이 모두 옛 줄이다. 마지막 기록은 20:57:17 UTC이고 release보다 앞이다 & 없음. 고칠 것이 없다 & 확인함: 요청 1번 뒤 줄 수 724 그대로, \code{page 200}. \textbf{20:57 UTC의 옛 timeout 원인은 모름} \\
\end{longtable}

\subsection{사례마다 가져갈 것}\label{sec:dbg-lessons}

\begin{longtable}{@{}L{0.5cm}L{14.0cm}@{}}
  \caption{열세 사례의 교훈. 열은 사례 번호와 교훈이다.}\label{tab:lessons}\\
  \toprule
  \# & 교훈 \\
  \midrule
  \endfirsthead
  \toprule
  \# & 교훈 \\
  \midrule
  \endhead
  \bottomrule
  \endfoot
  % src: books/debugging/01-...13-*.md "Lesson"
  1 & \code{/}와 \code{/robots.txt}를 먼저 비교한다. 그렇게 갈리면 한 번에 middleware를 가리킨다. \\
  \midrule
  2 & 페이지가 아니라 데이터 섬마다 가드를 둔다. 실패한 목록은 혼자 물러나야 한다. \\
  \midrule
  3 & ``페이지가 죽으면 안 된다''와 ``데이터가 틀렸다''를 나눈다. 앞의 것을 지금 고치고 뒤의 것은 따로 추적한다. \\
  \midrule
  4 & 서버를 건드리기 전에 토큰을 풀어 \code{iss}와 \code{azp}를 읽는다. 불일치가 거기에 적혀 있다. \\
  \midrule
  5 & origin을 포트까지 글자 그대로 읽는다. host가 맞아도 origin이 맞는 것은 아니다. \\
  \midrule
  6 & host 이름이 bind 주소를 정한다. 권한 오류가 주소를 말하면 주소를 직접 지정한다. \\
  \midrule
  7 & 새 host 이름은 세 곳을 건드린다: hosts 파일, 인증서, dev-origin 목록. \\
  \midrule
  8 & \code{pm2 env}로 실행 중인 값을 확인한다. 고친 파일은 프로세스에 대해 아무것도 증명하지 않는다. \\
  \midrule
  9 & 첫 \code{ZodError} 줄을 읽는다. 키 이름이 거기에 있다. 부팅 반복은 증명되기 전까지 설정 결함이다. \\
  \midrule
  10 & 인수 여럿은 배열에 담고 \texttt{"\$\{Z[@]\}"}로 푼다. zsh와 bash 모두에서 된다. \\
  \midrule
  11 & 로그인은 되는데 갱신이 실패하면 꺼진 API를 의심한다. 먼저 \code{gcloud services list}를 본다. \\
  \midrule
  12 & monorepo에서는 패키지 바이너리를 그 workspace로 돌린다. root에 있다고 가정하지 않는다. \\
  \midrule
  13 & release 뒤 pm2 err 로그에 \code{ETIMEDOUT}이 보이지만 오류는 옛 것이다. 시각 없는 로그 꼬리는 지금 실패한다는 증거가 아니다. 파일 수정 시각과, 요청 한 번 앞뒤의 줄 수를 먼저 본다. 길을 나눠 hop마다 \code{curl}을 하나씩 쓴다. \code{curl -6}이 약 1 ms에 실패하면 ``느림''이 아니라 ``경로 없음''이다. \\
\end{longtable}

\subsection{사례 13: release 뒤 오류 로그의 \texttt{ETIMEDOUT}}\label{sec:dbg-case13}

% src: scratchpad web-release-2026-10-13.facts.md §4; books/debugging/13-stale-error-log-etimedout.md
\paragraph{결과.}
고칠 것이 없었다.
로그의 724줄이 모두 옛 줄이었고, release 뒤 새 오류는 0줄이었다.
다만 20:57 UTC의 옛 timeout이 왜 났는지는 모른다.

\paragraph{무엇을 알려고 하나.}
release 뒤 오류 로그의 \texttt{ETIMEDOUT}이 지금 나는 오류인가, 옛 오류인가.

\paragraph{증상.}
\texttt{pm2 logs blog-web --err --lines 40 --nostream}이 아래 줄을 보였다.
파일은 \texttt{\textasciitilde/.pm2/logs/blog-web-error-3.log}다.
줄 앞의 \texttt{x} 기호는 생략했다.

\begin{lstlisting}[language={}]
[RecentPostsList] feed fetch failed: fetch failed
[TypeError: fetch failed] { digest: '3227098399', [cause]: AggregateError: { code: 'ETIMEDOUT', [errors]: [ [Error], [Error], [Error], [Error] ] } }
Error [TimeoutError]: The operation was aborted due to timeout { digest: '3827087268', code: 23, ... }
\end{lstlisting}

\paragraph{길.}
\Cref{fig:case13-path}가 길을 보인다.
\texttt{apps/dashboard\_client\_web/apis/http.server.ts:64}가 \texttt{fetch(\$\{API\_BASE\}\$\{path\})}를 부른다.
그래서 VM이 자기 API를 공개 도메인으로 부른다.
\texttt{AggregateError}의 오류 4개는 \texttt{blog-api.prettylog.com}이 풀린 주소 하나에 하나씩일 것이다. 아마 Cloudflare의 IPv4 주소 2개와 IPv6 주소 2개다. 추정이고 확인하지 않았다.

\begin{figure}[H]
  \centering
  \begin{tikzpicture}[node distance=0.4cm, font=\footnotesize]
    \node[flowbox, minimum width=2.3cm] (h1) {blog-web\\\scriptsize SSR (\texttt{http.server.ts:64})};
    \node[flowbox, minimum width=2.3cm, right=of h1] (h2) {\texttt{blog-api}\\\texttt{.prettylog.com}};
    \node[cfbox, minimum width=2.3cm, right=of h2] (h3) {Cloudflare};
    \node[flowbox, minimum width=2.3cm, right=of h3] (h4) {Tunnel};
    \node[vmbox, minimum width=2.3cm, right=of h4] (h5) {blog-api\\\scriptsize :4000};
    \foreach \x/\y in {h1/h2,h2/h3,h3/h4,h4/h5} \draw[flowarrow] (\x) -- (\y);
    \node[goodbox, minimum width=2.3cm, below=1.1cm of h1] (c1) {\texttt{page 200}\\\scriptsize 0.580770 s};
    \node[goodbox, minimum width=2.3cm, below=1.1cm of h2] (c2) {\texttt{v4 200}\\\scriptsize 0.201960 s};
    \node[badbox, minimum width=2.3cm, below=1.1cm of h3] (c3) {\texttt{v6 000}\\\scriptsize 0.001322 s};
    \node[goodbox, minimum width=2.3cm, below=1.1cm of h5] (c5) {\texttt{local 200}};
    \draw[flowarrow, dashed] (c1) -- (h1);
    \draw[flowarrow, dashed] (c2) -- (h2);
    \draw[flowarrow, dashed] (c3) -- (h2);
    \draw[flowarrow, dashed] (c5) -- (h5);
  \end{tikzpicture}
  \caption{사례 13의 길과 hop마다 하나씩 쓴 \texttt{curl}. 초록은 통과, 빨강은 실패다.}
  \label{fig:case13-path}
  % src: scratchpad web-release-2026-10-13.facts.md §4
\end{figure}

\paragraph{함정.}
\texttt{pm2 logs --lines 40}은 파일의 마지막 40줄을 시각 없이 보인다.
이 출력으로는 옛 오류와 새 오류를 가를 수 없다.

\paragraph{확인.}
오너가 VM에서 아래를 돌렸다.
\Cref{tab:case13-out}에 명령과 실제 출력을 나란히 적었다.
표의 \texttt{L}은 앞의 로그 파일 경로다.

\begin{table}[H]
  \caption{사례 13의 확인 명령과 실제 출력(2026/10/13).}
  \label{tab:case13-out}
  \centering\footnotesize
  \begin{tabularx}{\linewidth}{@{}L{7.2cm}Y@{}}
    \toprule
    명령 & 출력 \\
    \midrule
    \code{ls -l --time-style=full-iso $L} & \code{2026-10-13 20:57:17.958 +0000}, 20646 bytes \\
    \midrule
    \code{wc -l $L} & 724 \\
    \midrule
    \code{curl -m 20 ... http://localhost:3000/} & \code{page 200 0.580770s} \\
    \midrule
    \code{wc -l $L} (요청 뒤) & 724 (그대로) \\
    \midrule
    \code{curl -4 -m 10 ... https://blog-api.prettylog.com/ping} & \code{v4 200 0.201960s} \\
    \midrule
    \code{curl -6 -m 10 ... https://blog-api.prettylog.com/ping} & \code{v6 000 0.001322s} \\
    \midrule
    \code{curl -m 5 ... http://127.0.0.1:4000/ping} & \code{local 200} \\
    \midrule
    \code{pm2 ls} & \code{blog-api} online, uptime 2h, 재시작 횟수 7. \code{blog-web} online, uptime 6m, 재시작 횟수 8, 154.7mb \\
    \bottomrule
  \end{tabularx}
\end{table}

\paragraph{읽는 법.}
\Cref{tab:case13-read}는 결과마다 뜻을 한 줄로 적는다.

\begin{table}[H]
  \caption{결과와 뜻. 이 표를 같은 증상을 만날 때 쓴다.}
  \label{tab:case13-read}
  \centering\small
  \begin{tabularx}{\linewidth}{@{}L{5.4cm}Y@{}}
    \toprule
    결과 & 뜻 \\
    \midrule
    줄 수가 그대로이고 \code{page 200} & 옛 오류다. 이 release는 정상이다 \\
    \midrule
    \code{v4 200}, \code{v6 000} & IPv6 경로가 없다. 보통은 v4로 되돌아간다 \\
    \midrule
    \code{v4 000}, \code{local 200} & API는 살아 있다. VM에서 Cloudflare로 나가는 길이 막혔다 \\
    \midrule
    \code{local 000} & \code{blog-api}가 죽었다. \code{pm2 logs blog-api --err}를 읽는다 \\
    \midrule
    가끔만 실패 & 일시적이다. \#70의 429 시간대와 견준다 \\
    \bottomrule
  \end{tabularx}
\end{table}

\paragraph{판정.}
\begin{itemize}[nosep]
  \item 724줄은 모두 옛 줄이다.
  \item 마지막 기록 20:57:17 UTC는 \commit{4011c84} 빌드(run 생성 23:19:45 UTC)보다 앞이다. 곧 release보다 앞이다.
  \item 새 요청 한 번이 0줄을 더했다.
  \item \code{v6 000}이 0.001 s 만에 났다. timeout이 아니라 즉시 실패다. GCP VM에 IPv6 경로가 없다.
  \item v4, 로컬 API, SSR이 모두 동작한다. 고칠 것이 없다.
\end{itemize}

\begin{callout}[모르는 것]
20:57 UTC의 옛 timeout이 왜 났는지는 정하지 못했다.
추측으로 채우지 않는다.
\end{callout}

\paragraph{교훈.}
\begin{enumerate}[nosep]
  \item 시각 없는 로그 꼬리는 지금 실패한다는 증거가 아니다. 파일 수정 시각과, 요청 한 번 앞뒤의 줄 수를 본다.
  \item 길을 나눈다: SSR 페이지, 공개 API의 v4와 v6, 로컬 API. hop마다 timeout을 건 \texttt{curl}을 하나씩 쓴다.
  \item \texttt{curl -6}이 약 1 ms에 실패하면 ``경로 없음''이다. ``느림''이 아니다.
\end{enumerate}

% src: scratchpad logging-views-errorjson-2026-10-13.facts.md §1 §2
\paragraph{후속.}
PM2 \texttt{log\_date\_format}은 쓰지 않았다.
모든 줄에 날짜가 붙어 Ops Agent의 \texttt{parse\_json}이 JSON 줄에서 실패하기 때문이다.
대신 커밋 \commit{35b3630}이 시각 없는 오류 줄을 시각 있는 JSON 줄로 바꿨다.
VM에 배포하기 전이라 새 줄은 아직 관찰하지 않았다.
배포 뒤에는 \texttt{ls -l}과 \texttt{wc -l} 우회 대신 아래 질의를 쓴다.
빈 결과는 지난 1시간에 오류가 없다는 뜻이고, release보다 앞선 시각은 옛 오류다.
자세한 내용은 \Cref{sec:06-errorjson}에 있다.

\begin{lstlisting}[language={}]
gcloud logging read 'LOG_ID("blog-web") AND severity>=WARNING' --freshness=1h --limit 20 \
  --format='value(timestamp, jsonPayload.message, jsonPayload.cause.code)'
\end{lstlisting}
